\DocumentMetadata{
  pdfversion=2.0,
  pdfstandard=ua-2,
  lang=en-US,
  tagging=on,
  tagging-setup={math/setup=mathml-SE,tabsorder=structure}
}
\documentclass[11pt,letterpaper]{article}
\usepackage[margin=0.68in,includefoot]{geometry}
\usepackage{newcomputermodern}
\usepackage{amsmath,amscd,mathtools}
\usepackage{tikz,adjustbox}
\providecommand{\square}{\mdlgwhtsquare}
\usepackage[hidelinks]{hyperref}
\hypersetup{
  pdftitle={Numerical Analysis PhD Exam, May 2013},
  pdfauthor={University of Florida Department of Mathematics},
  pdfsubject={Graduate mathematics examination},
  pdfdisplaydoctitle=true,
  bookmarksopen=true
}
\setcounter{secnumdepth}{0}
\setlength{\parindent}{0pt}
\setlength{\parskip}{0.28em}
\setlength{\emergencystretch}{3em}
\setlength{\leftmargini}{2.15em}
\setlength{\leftmarginii}{2.65em}
\setlength{\leftmarginiii}{3em}
\renewcommand{\labelenumi}{\textbf{\theenumi.}}
\pagestyle{empty}
\raggedbottom
\allowdisplaybreaks
\newenvironment{examproblems}[1]{%
  \begin{enumerate}%
  \setcounter{enumi}{#1}%
  \setlength{\topsep}{0.35em}%
  \setlength{\partopsep}{0pt}%
  \setlength{\itemsep}{0.48em}%
  \setlength{\parsep}{0.18em}%
}{\end{enumerate}}
\newenvironment{examsubparts}{%
  \begin{enumerate}%
  \setlength{\topsep}{0.18em}%
  \setlength{\partopsep}{0pt}%
  \setlength{\itemsep}{0.16em}%
  \setlength{\parsep}{0.1em}%
}{\end{enumerate}}
\newenvironment{examinstructions}{%
  \par\begingroup\itshape
}{%
  \par\endgroup\vspace{0.18em}
}
\makeatletter
\renewcommand\section{\@startsection{section}{1}{0pt}%
  {-1.25ex plus -.25ex minus -.1ex}%
  {0.55ex plus .1ex}%
  {\normalfont\large\bfseries}}
\renewcommand{\@maketitle}{%
  \newpage\null\vskip -1em%
  {\footnotesize\bfseries
    \href{https://www.ufl.edu/}{University of Florida}, \href{https://math.ufl.edu/}{Department of Mathematics}\par}%
  \vskip -0.5em%
  \tagstructbegin{tag=Title}%
    {\LARGE\bfseries\href{https://gma.math.ufl.edu/exams/numerical-analysis-phd/}{Numerical Analysis PhD Exam}, May 2013\par}%
  \tagstructend%
  \vskip 0.25em%
  \tagmcbegin{artifact}%
    \hrule height 1pt%
  \tagmcend%
  \vskip 1em%
}
\makeatother
\title{Numerical Analysis PhD Exam, May 2013}
\author{}
\date{}
\begin{document}
\maketitle
\thispagestyle{empty}
\begin{examinstructions}
Do 5 of 7
\end{examinstructions}
\begin{examproblems}{0}
\item
(General Fixed Point Theory) Assume that \(g(x)\) is continuously differentiable on \([a,b]\), that \(g([a,b])\subseteq [a,b]\), and that 
\[
\lambda=\max_{a\leq x\leq b}|g'(x)|<1.
\]
 Show that the following are true:
\begin{examsubparts}
\item[\textnormal{(i)}]
\(x=g(x)\) has a unique solution~\(\alpha\) in \([a,b]\),
\item[\textnormal{(ii)}]
For any initial choice \(x_0\) in \([a,b]\), with \(x_{n+1}=g(x_n)\), \(\lim_{n\to\infty}x_n=\alpha\), and
\item[\textnormal{(iii)}]
\[
\lim_{n\to\infty}\frac{\alpha-x_{n+1}}{\alpha-x_n}=g'(\alpha).
\]
\item[\textnormal{(iv)}]
Does the statement remain true if the interval \((a,b)\) is open?
\item[\textnormal{(v)}]
Assume in addition that \(g'(\alpha)=0\) and show that the sequence converges quadratically.
\item[\textnormal{(vi)}]
(Newton’s Method) Assume that \(f(x)\), \(f'(x)\), and \(f''(x)\) are continuous in \([a,b]\), and that for some \(\alpha\in[a,b]\), \(f(\alpha)=0\), and \(f'(\alpha)\neq0\). Then if \(x_0\) is chosen close enough to~\(\alpha\), the iterates 
\[
x_{n+1}=x_n-\frac{f(x_n)}{f'(x_n)}
\]
 converge to~\(\alpha\). Moreover they converge quadratically.
\item[\textnormal{(vii)}]
State Newton’s method for \(f(x)=0\) if \(f:\mathbb{R}^n\to\mathbb{R}^n\).
\end{examsubparts}
\item
Theorem: (Lagrange Error Formula) Let \(x_0,x_1,x_2,\ldots,x_n\) be distinct real numbers, \(l_j(x)\) be the corresponding Lagrange polynomials and suppose that~\(f\) is a given real-valued function with \(n+1\) continuous derivatives. Let~\(I\) be an interval containing all of the \(x_k\), and~\(t\). Then there is a \(\xi\in I\) such that 
\[
f(t)-\sum_{j=0}^{n}f(x_j)l_j(t)=\frac{(t-x_0)(t-x_1)\cdots(t-x_n)}{(n+1)!}f^{(n+1)}(\xi).
\]
\item
Let \(\{p_n\}\) be an orthogonal family on \([a,b]\) constructed by using the Gram–Schmidt process on \(1,t,t^2,t^3,\ldots\). Prove that all of the zeros of \(p_n(t)\) are contained in \([a,b]\).
\item
Given Simpson’s rule for numerical integration, i.e. 
\[
\int_{t_n}^{t_{n+2}}f(x)\,dx\approx\frac{2h}{6}\bigl(f(t_n)+4f(t_{n+1})+f(t_{n+2})\bigr),\qquad t_n=t_0+nh,
\]
\begin{examsubparts}
\item[\textnormal{(i)}]
Explain where the formula comes from (you don’t need to derive it).
\item[\textnormal{(ii)}]
Prove that it is exact for cubic polynomials.
\item[\textnormal{(iii)}]
Show that the local error is \(O(h^5)\).
\end{examsubparts}
\item
Gaussian Quadrature: Show that you can find~\(n\) points on an interval \([a,b]\) (tell us which points), and a formula which uses only the value of the function \(f(x)\) at these points, and weights \(w_k\) such that the approximation formula 
\[
\int_a^b f(x)\,dx\approx I_n(f)\equiv\sum_{k=1}^{n}w_kf(x_k)
\]
 is exact for polynomials of order \(2n-1\).
\item
Assume that you are solving the initial value problem \(y'(t)=f(t,y)\), with \(y(0)\) given. Assume further that \(f(t,y)\) satisfies the Lipschitz condition \(|f(t,y_1)-f(t,y_2)|\leq K|y_1-y_2|\) for all \(t\in[a,b]\). Explain Euler’s method and derive a cumulative error formula, not just a local error formula.
\item
(General Multistep Methods) Assume that you are solving the initial value problem \(y'(t)=f(t,y)\), with \(y(0)\) given. Consider a general formula of the type 
\[
y_{n+1}=\sum_{j=0}^{p}a_jy_{n-j}+h\sum_{j=-1}^{p}b_jf(t_{n-j},y_{n-j}).
\]
 Furthermore, let us define the local truncation error as 
\[
T_n(Y)=Y(t_{n+1})-\left(\sum_{j=0}^{p}a_jY(t_{n-j})+h\sum_{j=-1}^{p}b_jf(t_{n-j},Y(t_{n-j}))\right),
\]
 where \(Y(t_n)\) is the exact value of the initial value problem and \(y_n\) is the approximated value of the problem at \(t_n\). We let 
\[
\tau_n(Y)=\frac{1}{h}T_n(Y).
\]
 Given this prove the following:

\par
Let \(m\geq1\) be a given integer. In order that \(\max|\tau(Y)|\to0\) as \(h\to0\) for all continuously differentiable \(Y(x)\), it is necessary and sufficient that 
\[
\sum_{j=0}^{p}a_j=1,\qquad\text{and}\qquad-\sum_{j=0}^{p}ja_j+\sum_{j=-1}^{p}b_j=1. \tag{1}
\]
 Furthermore, for \(\tau(h)=O(h^m)\) for functions \(Y(x)\) that are \(m+1\) times continuously differentiable, it is necessary and sufficient that (1) hold and 
\[
\sum_{j=0}^{p}(-j)^ia_j+i\sum_{j=-1}^{p}(-j)^{i-1}b_j=1,\qquad \text{for }i=2,3,\ldots,m.
\]
\end{examproblems}
\begin{examinstructions}
Appendix: Recall that the Hermite polynomials can be written as 
\[
H_n(x)=\sum_{j=1}^{n}f(x_j)h_j(x)+\sum_{j=1}^{n}f'(x_j)\tilde h_j(x).
\]
 If the points \(x_j\) are chosen to be the zeros of the orthogonal polynomials on \([a,b]\), then 
\[
\tilde h_j(x)=\frac{\psi_n(x)l_j(x)}{\psi_n'(x_j)},
\]
 where \(l_j(x)\) is the Lagrange interpolant for \(x_j\).
\end{examinstructions}
\end{document}
